\documentclass[11pt,a4paper]{article}
\usepackage[utf8]{inputenc}
\usepackage[T1]{fontenc}
\usepackage{geometry}
\usepackage{booktabs}
\usepackage{longtable}
\usepackage{float}
\usepackage{array}
\usepackage{amsmath}
\usepackage{graphicx}
\usepackage{hyperref}
\usepackage{xcolor}
\geometry{margin=1in}

\title{Custom-WikiCS Graph Analysis -- Part 7\\[4pt]
\large Revised Undirected Structural Evaluation: 10-Model Comparison with Recommendation Metrics}
\author{Sahin Uregil}
\date{\today}

\begin{document}
\maketitle

\section{Introduction}
This report presents a revised evaluation branch for the Custom-WikiCS
recommendation project. The earlier comparison reports were useful, but they
mixed a directed edge-holdout setup with structural models that ultimately
operate on undirected graph structure. In particular, Node2Vec was trained on an
undirected graph, while the old train/test split was defined over directed
outgoing edges. That created a methodological mismatch and made it difficult to
separate true structural generalization from artifacts of the split design.

The new branch therefore changes the evaluation setup in three ways:

\begin{itemize}
    \item The structural branch now uses an \textbf{undirected train/test split}
          built directly on the undirected graph.
    \item The Node2Vec family is retrained on that new branch, and both
          \textbf{uncorrected} and \textbf{power-law corrected} versions are
          reevaluated.
    \item The report adds a second table of \textbf{recommendation-quality
          metrics} so that models are evaluated not only as edge-recovery
          systems, but also as recommendation systems.
\end{itemize}

This revised setup also introduces two new pooled hybrid models that balance
BGE-M3 and Node2Vec dimensionally before fusion.

\section{Revised Evaluation Setup}

\subsection{Undirected Structural Split}
The full structural graph is first converted to an undirected graph. Undirected
edges are then split into train and test sets with an approximately 80/20
ratio. To avoid weakening the structural train graph unnecessarily, the split
preserves train connectivity where possible and repairs any accidental new train
isolates.

The resulting branch statistics are:

\begin{itemize}
    \item Full undirected structural edges: 216,123
    \item Train edges: 172,900
    \item Test edges: 43,223
    \item Expanded evaluation pairs: 86,446
    \item Evaluation nodes: 9,022
\end{itemize}

For ranking-based evaluation, each undirected test edge $\{u,v\}$ is expanded
into two directional queries: $u \rightarrow v$ and $v \rightarrow u$.

\subsection{Power-Law Correction}
The power-law correction parameter is refitted on the new undirected train
graph. The fitted value for this branch is:

\[
\alpha_{\text{structural}} \approx 2.9966
\]

In this report, the terms ``corrected'' and ``uncorrected'' refer only to
whether this power-law penalty is applied to the Node2Vec representation.

\subsection{Pooled Hybrid Models}
The previous combined representation used:

\begin{itemize}
    \item BGE-M3: 1024 dimensions
    \item Node2Vec: 128 dimensions
\end{itemize}

To test whether this dimensional imbalance was hurting hybrid fusion, a pooled
variant was introduced. BGE-M3 is compressed from 1024 to 128 dimensions using
simple mean-pooling over contiguous blocks, then concatenated with 128-dim
Node2Vec. This creates a balanced 256-dimensional hybrid representation.

\section{Models Compared}

\begin{table}[H]
\centering
\caption{Summary of the 10 models evaluated in the revised branch.}
\begin{tabular}{>{\raggedright\arraybackslash}p{0.05\textwidth}
                >{\raggedright\arraybackslash}p{0.34\textwidth}
                >{\raggedright\arraybackslash}p{0.16\textwidth}
                >{\raggedright\arraybackslash}p{0.35\textwidth}}
\toprule
\textbf{\#} & \textbf{Model} & \textbf{Category} & \textbf{Description} \\
\midrule
1 & BGE-M3 Only & Semantic & Cosine similarity on 1024-dim BGE-M3 text embeddings \\
2 & Node2Vec (uncorrected) & Structural & Undirected-branch Node2Vec without power-law penalty \\
3 & Node2Vec (corrected) & Structural & Undirected-branch Node2Vec with power-law penalty \\
4 & BGE-M3 + Node2Vec (uncorrected) & Combined & 1024+128 concatenation on the new undirected branch \\
5 & BGE-M3 + Node2Vec (corrected) & Combined & 1024+128 concatenation with corrected Node2Vec component \\
6 & GCN Only & GNN & Structural GCN retrained on the revised branch \\
7 & GCN + BGE-M3 (rank fusion) & Fusion & Reciprocal Rank Fusion between revised-branch GCN and BGE-M3 \\
8 & GraphSAGE Only & GNN & Structural GraphSAGE retrained on the revised branch \\
9 & BGE pooled + Node2Vec (uncorrected) & Combined & 128-dim pooled BGE + 128-dim uncorrected Node2Vec \\
10 & BGE pooled + Node2Vec (corrected) & Combined & 128-dim pooled BGE + 128-dim corrected Node2Vec \\
\bottomrule
\end{tabular}
\end{table}

\section{Metric Families}

\subsection{Existing Ranking Metrics}
For continuity with the earlier report series, the following metrics are
retained:

\begin{itemize}
    \item Hit@1
    \item Hit@10
    \item Hit@50
    \item Hit@100
    \item MRR
    \item NDCG
    \item Recall
    \item MAP
\end{itemize}

\subsection{New Recommendation Metrics}
To better reflect the actual recommendation objective, the following metrics are
added:

\begin{itemize}
    \item Popularity Lift
    \item Gini Index
    \item Aggregate Diversity
    \item Novelty
    \item Intra-List Distance
\end{itemize}

\section{10-Model Results}

\subsection{Combined Metric Table}

\begin{table}[H]
\centering
\caption{Combined old and new metrics on the revised undirected structural branch. Ranking metrics use 86,446 expanded evaluation pairs; recommendation metrics use top-20 lists over 9,022 evaluation nodes.}
\setlength{\tabcolsep}{3pt}
\renewcommand{\arraystretch}{1.12}
\resizebox{\textwidth}{!}{
\begin{tabular}{>{\raggedright\arraybackslash}p{0.28\textwidth}ccccccccccccc}
\toprule
\textbf{Model} &
\rotatebox{45}{\textbf{H@1}} &
\rotatebox{45}{\textbf{H@10}} &
\rotatebox{45}{\textbf{H@50}} &
\rotatebox{45}{\textbf{H@100}} &
\rotatebox{45}{\textbf{MRR}} &
\rotatebox{45}{\textbf{NDCG}} &
\rotatebox{45}{\textbf{Recall}} &
\rotatebox{45}{\textbf{MAP}} &
\rotatebox{45}{\textbf{Pop. Lift}} &
\rotatebox{45}{\textbf{Gini}} &
\rotatebox{45}{\textbf{Agg. Div.}} &
\rotatebox{45}{\textbf{Novelty}} &
\rotatebox{45}{\textbf{ILD}} \\
\midrule
BGE pooled + Node2Vec (corrected) & 1.03\% & 8.67\% & 31.46\% & 47.81\% & 0.0410 & 0.1160 & 0.4781 & 0.0392 & 1.3301 & 0.4110 & 11,326 & 14.2738 & 0.4372 \\
BGE-M3 + Node2Vec (corrected) & 1.07\% & 8.67\% & 32.41\% & 49.81\% & 0.0421 & 0.1202 & 0.4981 & 0.0403 & 1.3081 & 0.4127 & 11,268 & 14.2783 & 0.4316 \\
BGE pooled + Node2Vec (uncorrected) & 0.76\% & 7.83\% & 32.64\% & 50.74\% & 0.0376 & 0.1176 & 0.5074 & 0.0358 & 1.3397 & 0.4613 & 11,186 & 14.0351 & 0.4558 \\
BGE-M3 + Node2Vec (uncorrected) & 0.76\% & 7.83\% & 32.64\% & 50.74\% & 0.0376 & 0.1176 & 0.5074 & 0.0358 & 1.3400 & 0.4605 & 11,164 & 14.0340 & 0.4557 \\
Node2Vec (uncorrected) & 0.75\% & 7.81\% & 32.52\% & 50.60\% & 0.0375 & 0.1172 & 0.5060 & 0.0357 & 1.3374 & 0.4482 & 10,966 & 14.0466 & 0.4569 \\
Node2Vec (corrected) & 0.75\% & 7.81\% & 32.52\% & 50.60\% & 0.0375 & 0.1172 & 0.5060 & 0.0356 & 1.3374 & 0.4482 & 10,966 & 14.0466 & 0.4569 \\
GCN + BGE-M3 (rank fusion) & 0.85\% & 7.63\% & 23.14\% & 33.77\% & 0.0335 & 0.0862 & 0.3377 & 0.0317 & 1.6650 & 0.4228 & 11,655 & 13.9222 & 0.4366 \\
BGE-M3 Only & 1.35\% & 7.57\% & 19.69\% & 27.73\% & 0.0366 & 0.0788 & 0.2773 & 0.0351 & 1.9413 & 0.5172 & 10,905 & 13.8459 & 0.3713 \\
GCN Only & 0.38\% & 3.46\% & 14.12\% & 22.91\% & 0.0180 & 0.0530 & 0.2291 & 0.0163 & 1.1879 & 0.2607 & 11,685 & 14.3293 & 0.5135 \\
GraphSAGE Only & 0.17\% & 1.22\% & 4.96\% & 8.72\% & 0.0076 & 0.0199 & 0.0872 & 0.0061 & 1.3134 & 0.2517 & 11,628 & 14.2087 & 0.5313 \\
\bottomrule
\end{tabular}
}
\end{table}

\subsection{Per-Metric Rank Summary}

To make the multi-metric comparison easier to scan, Table~\ref{tab:rank-summary}
converts each raw metric into a per-column ranking. Rank 1 is best in each
metric column. For Hit@K, MRR, NDCG, Recall, MAP, Aggregate Diversity, Novelty,
and Intra-List Distance, higher values receive better ranks. For Popularity
Lift and Gini, lower values receive better ranks. The \textbf{Total Rank}
column reports the final ordering by the summed per-metric ranks, so lower is
better overall.

\begin{table}[H]
\centering
\caption{Per-metric rank summary across old and new evaluation metrics. Lower total rank is better overall.}
\label{tab:rank-summary}
\setlength{\tabcolsep}{3pt}
\renewcommand{\arraystretch}{1.15}
\resizebox{\textwidth}{!}{
\begin{tabular}{>{\raggedright\arraybackslash}p{0.28\textwidth}cccccccccccccc}
\toprule
\textbf{Model} &
\rotatebox{45}{\textbf{Total Rank}} &
\rotatebox{45}{\textbf{Hit@1}} &
\rotatebox{45}{\textbf{Hit@10}} &
\rotatebox{45}{\textbf{Hit@50}} &
\rotatebox{45}{\textbf{Hit@100}} &
\rotatebox{45}{\textbf{MRR}} &
\rotatebox{45}{\textbf{NDCG}} &
\rotatebox{45}{\textbf{Recall}} &
\rotatebox{45}{\textbf{MAP}} &
\rotatebox{45}{\textbf{Agg. Div.}} &
\rotatebox{45}{\textbf{Novelty}} &
\rotatebox{45}{\textbf{ILD}} &
\rotatebox{45}{\textbf{Pop. Lift}} &
\rotatebox{45}{\textbf{Gini}} \\
\midrule
5. BGE-M3 + Node2Vec (corrected) & 1 & 2 & 2 & 5 & 5 & 1 & 1 & 5 & 1 & 5 & 2 & 9 & 2 & 4 \\
10. BGE pooled + Node2Vec (corrected) & 2 & 3 & 1 & 6 & 6 & 2 & 6 & 6 & 2 & 4 & 3 & 7 & 4 & 3 \\
9. BGE pooled + Node2Vec (uncorrected) & 3 & 6 & 3 & 1 & 1 & 4 & 3 & 1 & 4 & 6 & 7 & 5 & 7 & 9 \\
4. BGE-M3 + Node2Vec (uncorrected) & 4 & 5 & 4 & 2 & 2 & 3 & 2 & 2 & 3 & 7 & 8 & 6 & 8 & 8 \\
2. Node2Vec (uncorrected) & 5 & 7 & 5 & 3 & 3 & 5 & 4 & 3 & 5 & 8 & 5 & 3 & 5 & 6 \\
3. Node2Vec (corrected) & 6 & 8 & 6 & 3 & 3 & 6 & 5 & 3 & 6 & 8 & 5 & 3 & 5 & 6 \\
6. GCN Only & 7 & 9 & 9 & 9 & 9 & 9 & 9 & 9 & 9 & 1 & 1 & 2 & 1 & 2 \\
7. GCN + BGE-M3 (rank fusion) & 8 & 4 & 7 & 7 & 7 & 8 & 7 & 7 & 8 & 2 & 9 & 8 & 9 & 5 \\
8. GraphSAGE Only & 9 & 10 & 10 & 10 & 10 & 10 & 10 & 10 & 10 & 3 & 4 & 1 & 3 & 1 \\
1. BGE-M3 Only & 10 & 1 & 8 & 8 & 8 & 7 & 8 & 8 & 7 & 10 & 10 & 10 & 10 & 10 \\
%
\bottomrule
\end{tabular}
}
\end{table}

\section{What Changed Under the New Setup}

\subsection{Node2Vec No Longer Dominates}
The most important change is that \textbf{pure Node2Vec is no longer the clear
winner}. Under the older branch, uncorrected Node2Vec was the strongest model by
a large margin. Under the revised undirected structural branch, both Node2Vec
variants settle around:

\begin{itemize}
    \item Hit@10 $\approx 7.81\%$
    \item MRR $\approx 0.0375$
\end{itemize}

This is no longer a dominant result. The revised branch therefore weakens the
earlier conclusion that structural Node2Vec alone is the strongest model in the
project.

\subsection{Hybrid Models Now Form the Top Tier}
The strongest models are now the \textbf{corrected hybrid models}:

\begin{itemize}
    \item \textbf{BGE pooled + Node2Vec (corrected)} achieves the best Hit@10:
          $8.67\%$
    \item \textbf{BGE-M3 + Node2Vec (corrected)} is essentially tied on Hit@10
          and slightly stronger on MRR, Hit@50, Hit@100, Recall and MAP
\end{itemize}

This is a much more plausible outcome for the intended project direction. It
suggests that the best-performing family is not pure structure alone, but
hybrid semantic-structural retrieval.

\subsection{Pooling Helps Slightly, But Not Universally}
The pooling experiment was worthwhile. Dimensionally balancing BGE-M3 and
Node2Vec does improve the top-line Hit@10 very slightly:

\begin{itemize}
    \item pooled corrected hybrid: Hit@10 = $8.6736\%$
    \item plain corrected hybrid: Hit@10 = $8.6690\%$
\end{itemize}

However, the gain is very small, and the plain corrected hybrid is still
slightly stronger on several secondary ranking metrics. So the correct
interpretation is:

\begin{itemize}
    \item pooling helps a little
    \item but it does not completely replace the non-pooled corrected hybrid
\end{itemize}

\subsection{Degree Correction Behaves Differently Now}
Another important shift is that \textbf{degree correction no longer severely
damages pure Node2Vec}. On this new branch, corrected and uncorrected Node2Vec
are almost identical.

That means the earlier strong claim
``power-law correction hurts pure Node2Vec badly'' does not survive the revised
evaluation setup. The effect appears to have been highly dependent on the old
branch.

\section{Interpretation of the New Recommendation Metrics}

\subsection{GCN Has the Best Diversity Profile}
Although GCN Only is weak on edge-recovery metrics, it looks strong under the
new recommendation-oriented metrics:

\begin{itemize}
    \item lowest popularity lift among strong baselines
    \item highest aggregate diversity
    \item high novelty
    \item strong intra-list distance
\end{itemize}

So GCN appears to trade off ranking accuracy for broader and less popularity-led
recommendation behavior.

\subsection{BGE-M3 Is the Most Popularity-Driven}
BGE-M3 Only shows:

\begin{itemize}
    \item the highest popularity lift ($1.9413$)
    \item the highest Gini concentration ($0.5172$)
    \item the lowest novelty among the major families
    \item the lowest intra-list diversity ($0.3713$)
\end{itemize}

This suggests that pure semantic retrieval tends to collapse more strongly onto
popular or obvious recommendations.

\subsection{Hybrids Occupy the Best Middle Ground}
The corrected hybrid models appear to offer the best compromise:

\begin{itemize}
    \item clearly better ranking accuracy than GCN and GraphSAGE
    \item much less popularity bias than pure BGE-M3
    \item strong diversity and novelty without sacrificing top-line performance
\end{itemize}

This is exactly the kind of tradeoff that matters for the actual product
objective of article recommendation and learning-path generation.

\section{GNN Branch Notes}
The revised-branch GNN retraining succeeded and produced:

\begin{itemize}
    \item GCN Only test AUC $\approx 0.9304$
    \item GraphSAGE Only test AUC $\approx 0.9190$
\end{itemize}

This preserves the earlier pattern that GCN outperforms GraphSAGE on this
transductive structural task, even though both remain weaker than the top hybrid
models under ranking metrics.

\section{Conclusions}
The revised undirected structural branch changes the project story in an
important way.

\begin{itemize}
    \item \textbf{Pure Node2Vec is no longer the dominant model.}
    \item \textbf{Corrected hybrid models now form the strongest overall tier.}
    \item \textbf{Pooling provides a small but real improvement on Hit@10.}
    \item \textbf{The old claim that degree correction strongly harms Node2Vec
          does not appear robust.}
    \item \textbf{Recommendation-quality metrics reveal a second axis of model
          quality beyond edge recovery.}
\end{itemize}

The strongest overall interpretation of this report is therefore:

\begin{quote}
Under the revised, recommendation-oriented structural setup, the best models are
hybrid semantic-structural models, not pure structure alone. The corrected
hybrid family is now the most convincing direction for the project's final
system.
\end{quote}

Among those, the two main candidates are:

\begin{itemize}
    \item \textbf{BGE pooled + Node2Vec (corrected)} if Hit@10 is prioritized
    \item \textbf{BGE-M3 + Node2Vec (corrected)} if the slightly stronger MRR,
          Recall and MAP profile is prioritized
\end{itemize}

This revised branch is therefore both methodologically cleaner and better
aligned with the true article-recommendation objective of the project.

\end{document}
